\documentclass[
    language=english,
    doctype=course-material,
    institution=none,
]{../../omnilatex}

\RequirePackage{config/document-settings}

\begin{document}

\maketitle

\section{Course Description}

This course provides a rigorous introduction to the design, analysis, and
implementation of fundamental data structures and algorithms.  Topics include
arrays, linked lists, stacks, queues, trees, hash tables, graphs, sorting,
searching, and algorithmic complexity.  Students will implement data structures
in \textsf{Java} and analyze their asymptotic performance.

\section{Learning Outcomes}

By the end of this course, students will be able to:

\begin{itemize}
    \item Implement and use arrays, linked lists, stacks, queues, and hash tables.
    \item Analyze the time and space complexity of algorithms using Big-O, Big-$\Omega$,
          and Big-$\Theta$ notation.
    \item Apply tree structures (binary trees, BSTs, AVL trees) to solve search
          and retrieval problems.
    \item Implement graph algorithms including BFS, DFS, Dijkstra's shortest path,
          and topological sorting.
    \item Select the most appropriate data structure for a given computational problem.
    \item Evaluate trade-offs between time complexity, space complexity, and
          implementation difficulty.
\end{itemize}

\section{Required Materials}

\begin{itemize}
    \item Mark Allen Weiss, \emph{Data Structures and Algorithm Analysis in Java},
          4th~ed., Pearson, 2013.
    \item Robert Sedgewick and Kevin Wayne, \emph{Algorithms}, 4th~ed.,
          Addison-Wesley, 2011.
    \item Java Development Kit (JDK)~17 or later.
    \item IntelliJ IDEA or equivalent IDE (free educational licenses available).
\end{itemize}

\section{Grading Policy}

\begin{table}[htbp]
    \centering
    \caption{Grading breakdown}
    \begin{tabular}{@{}lc@{}}
        \toprule
        Component & Weight (\%) \\
        \midrule
        Programming Assignments (6) & 30 \\
        Quizzes (4) & 15 \\
        Midterm Examination & 20 \\
        Final Examination & 25 \\
        Participation \& Attendance & 10 \\
        \bottomrule
    \end{tabular}
\end{table}

\section{Course Policies}

\subsection{Academic Integrity}

All submitted work must represent your own effort.  You may discuss problems
at a high level but must write all code independently.  Use of AI coding
assistants must be disclosed; code submitted without disclosure will be
treated as a violation of academic integrity.

\subsection{Late Submissions}

Each student receives three late-day tokens (24 hours each) for the semester.
Tokens may not be applied to exams.  After tokens are exhausted, late work
receives a 10\% penalty per day.

\section{Weekly Schedule}

\small
\begin{tabular}{@{}clp{8cm}@{}}
    \toprule
    \textbf{Week} & \textbf{Dates} & \textbf{Topic} \\
    \midrule
    \scheduleentry{1}{Sep 1--3}{Course overview; complexity analysis: Big-O, Big-$\Omega$, Big-$\Theta$}
    \scheduleentry{2}{Sep 8--10}{Arrays and linked lists; dynamic arrays; amortized analysis}
    \scheduleentry{3}{Sep 15--17}{Stacks and queues; applications: balanced parentheses, BFS}
    \scheduleentry{4}{Sep 22--24}{Recursion; recurrence relations; the Master Theorem}
    \scheduleentry{5}{Sep 29--Oct 1}{\textbf{Quiz 1}; Trees: binary trees, traversals}
    \scheduleentry{6}{Oct 6--8}{Binary search trees; balanced BSTs (AVL, Red-Black)}
    \scheduleentry{7}{Oct 13--15}{Hash tables: chaining, open addressing, universal hashing}
    \scheduleentry{8}{Oct 20--22}{\textbf{Midterm Examination}}
    \scheduleentry{9}{Oct 27--29}{Heaps and priority queues; heap sort}
    \scheduleentry{10}{Nov 3--5}{Graphs: representations, BFS, DFS; \textbf{Quiz 2}}
    \scheduleentry{11}{Nov 10--12}{Shortest paths: Dijkstra's and Bellman--Ford algorithms}
    \scheduleentry{12}{Nov 17--19}{Minimum spanning trees: Kruskal's and Prim's algorithms}
    \scheduleentry{13}{Nov 24--26}{Sorting: merge sort, quicksort, lower bounds; \textbf{Quiz 3}}
    \scheduleentry{14}{Dec 1--3}{String matching: Knuth--Morris--Pratt, Rabin--Karp}
    \scheduleentry{15}{Dec 8--10}{Course review; \textbf{Quiz 4}; \textbf{Final exam preparation}}
    \midrule
    & \multicolumn{2}{l}{\textbf{Midterm:} October 22, in class} \\
    & \multicolumn{2}{l}{\textbf{Final:} December 15, 9:00--11:00\,AM} \\
    \bottomrule
\end{tabular}

\section{Weekly Topic Detail}

\subsection{Weeks 1--2: Foundations and Complexity}

The first two weeks establish the mathematical foundations for algorithm
analysis.  We will review asymptotic notation and introduce recurrence
relations.  Students will implement dynamic arrays and singly linked lists,
comparing their performance empirically against theoretical predictions.

\subsection{Weeks 3--4: Abstract Data Types and Recursion}

Stacks and queues are introduced as abstract data types with concrete
array-based and linked-list implementations.  Recursive problem-solving
techniques are formalized, leading to the Master Theorem for divide-and-conquer
recurrences~\cite{knuthFundamentalAlgorithms1973}.

\subsection{Weeks 5--7: Trees and Hash Tables}

Tree structures form the backbone of efficient search.  We cover binary search
trees, self-balancing variants, and hash tables.  Students will implement a
hash table with separate chaining and compare its performance to open-addressing
schemes.

\subsection{Weeks 9--12: Graphs and Advanced Algorithms}

Graph algorithms build on earlier topics.  BFS and DFS are used as subroutines
for shortest-path and minimum-spanning-tree algorithms.  We analyze Dijkstra's
algorithm using priority queues and prove its correctness.

\subsection{Weeks 13--15: Sorting, Strings, and Review}

We prove the $\Omega(n \log n)$ lower bound for comparison-based sorting and
implement efficient string-matching algorithms.  The final week is devoted to
comprehensive review and exam preparation.

\section{Supplementary References}

\begin{itemize}
    \item Thomas H.~Cormen et al., \emph{Introduction to Algorithms}, 4th~ed.,
          MIT Press, 2022.
    \item Jon Kleinberg and \'{E}va Tardos, \emph{Algorithm Design}, Pearson, 2005.
\end{itemize}

\printbibliography[heading=bibintoc,title={References}]

\end{document}
